%% =============================================================================
%% sm_d5_states.tex — SM-D, Subsection D.3: State-Level Poverty Coefficients
%% Body content only (\subsection level); parent structure in sm_d.tex
%%
%% =============================================================================

\subsection{State-level poverty coefficients}
\label{smd:states}

This subsection summarizes the state-level poverty coefficients from
model~M2 (block-diagonal SVC without policy moderators) for all
$S = 51$ states.  The extensive-margin coefficient
$\alpha_{\mathrm{pov},s}$ and intensive-margin coefficient
$\beta_{\mathrm{pov},s}$ are the state-specific effects of a
one-standard-deviation increase in community poverty on the log-odds of
IT participation and the logit of IT intensity, respectively.  The
cross-margin scatter plot in \cref{fig:cross-margin} visualizes the
joint distribution of these estimates.

Of the 51 states, 48 (94\%) have posterior-mean sign patterns consistent
with the classic reversal
($\hat{\alpha}_{\mathrm{pov},s} < 0$ and
$\hat{\beta}_{\mathrm{pov},s} > 0$); this near-universal pattern
reflects hierarchical shrinkage pulling small-state estimates toward
the negative population-average extensive-margin coefficient.
Among these, 15 (29\%) exhibit posterior reversal probabilities
exceeding 0.80, indicating strong evidence beyond the point-estimate
sign, and an additional 21 states have reversal probabilities between
0.60 and 0.80.

\paragraph{States with strongest reversal evidence.}
The five large-sample states with the highest reversal probabilities are
Texas ($N = 578$, $\Pr = 0.989$),
Illinois ($N = 471$, $\Pr = 0.970$),
California ($N = 1{,}110$, $\Pr = 0.937$),
Louisiana ($N = 77$, $\Pr = 0.923$), and
New Jersey ($N = 201$, $\Pr = 0.908$).
In these states, the extensive-margin coefficients
($\alpha_{\mathrm{pov},s} \in [-0.34, -0.13]$) are firmly negative
with 90\% credible intervals excluding zero, while the
intensive-margin coefficients
($\beta_{\mathrm{pov},s} \in [0.07, 0.08]$) are positive with
intervals barely including zero.

\paragraph{Notable exception: New York.}
New York ($N = 558$) stands out as the most prominent exception, with a
\emph{positive} extensive-margin coefficient
($\alpha_{\mathrm{pov}} = 0.101$; 90\% CI $[-0.051, 0.261]$,
$\Pr(\text{reversal}) = 0.143$), suggesting that higher poverty is
associated with \emph{greater} likelihood of serving infants and
toddlers---possibly reflecting New York's substantial public investment
in early childhood programs.  Arizona ($\Pr = 0.40$), Minnesota
($\Pr = 0.44$), and North Dakota ($\Pr = 0.48$) also show weak or
absent reversal patterns.

\paragraph{Small-state shrinkage.}
Small states (e.g., Maine, $N = 18$; Mississippi, $N = 17$; Montana,
$N = 22$) exhibit wider credible intervals and reversal probabilities
closer to the national average ($\approx 0.73$), reflecting the
expected shrinkage toward the hierarchical mean documented
in~\cref{app:survey}.  The full 51-state coefficient table is available
in the replication package at
\url{https://github.com/joonho112/hurdlebb-replication}.
